\chapter{transformaciones Coordenadas}

\section{Introducción y motivación}
% ====================================================================

Cuando estudiamos un sistema físico anhelamos que las ecuaciones del movimiento sean simples en su forma funcional, con el fin de poder resolverlas. Sin embargo, esto no siempre ocurre. No obstante, si cambiamos el sistema coordenado a otro, las ecuaciones podrían, en principio, simplificarse.

Por ejemplo, en el problema de fuerzas centrales conservativas la segunda ley de Newton se reduce a
\begin{equation}
m\vec a = f\!\left(\sqrt{x^2+y^2}\right)\hat x .
\end{equation}
Si usamos coordenadas cartesianas, las ecuaciones del movimiento se escriben de forma explícita como
\begin{equation}
m\ddot x = f\!\left(\sqrt{x^2+y^2}\right)x, \qquad
m\ddot y = f\!\left(\sqrt{x^2+y^2}\right)y, \qquad
z=0,
\end{equation}
lo que se traduce en dos ecuaciones diferenciales ordinarias acopladas y, en general, no lineales. Resolver estas ecuaciones no siempre es simple.

Si en cambio usamos coordenadas cilíndricas
\begin{equation}
x = r\cos\phi, \qquad y = r\sin\phi, \qquad z = 0, \qquad r\ge 0,\quad 0\le \phi < 2\pi,
\end{equation}
las ecuaciones del movimiento del problema central se simplifican a
\begin{equation}
m\!\left(\ddot r - r\dot\phi^{\,2}\right) = f(r), \qquad
m\!\left(r^2\ddot\phi + 2r\dot r\dot\phi\right) = 0 .
\end{equation}
La segunda ecuación se reduce a
\begin{equation}
\frac{\dd}{\dd t}\!\left(mr^2\dot\phi\right) = 0
\qquad\Longrightarrow\qquad
mr^2\dot\phi = l,
\end{equation}
es decir, usando coordenadas cilíndricas se obtiene una constante del movimiento, el momentum angular $l$; una de las ecuaciones del movimiento reduce así su orden diferencial de segundo a primer orden. Usando esta constante, las ecuaciones del movimiento en coordenadas cilíndricas se reducen finalmente a
\begin{equation}
m\ddot r = f(r) + \frac{l^2}{mr^3}, \qquad mr^2\dot\phi = l .
\end{equation}
Las ecuaciones del movimiento se simplifican bastante al pasar de coordenadas cartesianas a cilíndricas. Esto no significa que necesariamente se puedan resolver, pero la simplificación ayuda bastante.

Una \textbf{transformación de coordenadas} corresponde a un cambio de coordenadas que permite escribir las ecuaciones del movimiento en una forma más simple. Por ejemplo, al cambiar de coordenadas cartesianas a esféricas
\begin{equation}
x = r\cos\phi\sin\theta,\qquad y = r\sin\phi\sin\theta,\qquad z = r\cos\theta,
\qquad r\ge 0,\quad 0\le\phi<2\pi,\quad 0\le\theta\le\pi,
\end{equation}
se tiene el cambio $(x,y,z)\longrightarrow (r,\phi,\theta)$.

Una transformación de coordenadas no solo cambia los parámetros de las coordenadas, sino que también cambian las \emph{líneas coordenadas}. En el ejemplo anterior, al cambiar de cartesianas a esféricas, se pasa de tres rectas (los ejes $x,y,z$) a un rayo, una circunferencia y una media circunferencia.

% ====================================================================
\section{Transformación de coordenadas}
% ====================================================================

\subsection{Definición formal}

Formalmente, una transformación de coordenadas corresponde a un cambio de la forma
\begin{equation}
(x^1,x^2,x^3) \longleftrightarrow (x'^1,x'^2,x'^3),
\end{equation}
tal que
\begin{equation}
\left.
\begin{aligned}
x^1 &= x^1(x'^1,x'^2,x'^3) = x^1(x'^j) \\
x^2 &= x^2(x'^1,x'^2,x'^3) = x^2(x'^j) \\
x^3 &= x^3(x'^1,x'^2,x'^3) = x^3(x'^j)
\end{aligned}
\right\}
\quad\Longrightarrow\quad x^i = x^i(x'^j).
\end{equation}
El vector posición de un punto arbitrario cambia entonces en la forma
\begin{equation}
\rr(x^i) = \rr\big(x^i(x'^j)\big) = \rrp(x'^j)
\qquad\therefore\qquad \rr(x^i) = \rrp(x'^j).
\end{equation}
En general $\rr(x^i)\neq \rr(x'^j)$, puesto que la forma funcional del vector posición puede cambiar con las nuevas coordenadas $x'^i$. A menos que se diga explícitamente lo contrario, usaremos la nomenclatura
\begin{itemize}
\item $x^i$: coordenadas viejas,
\item $x'^i$: coordenadas nuevas.
\end{itemize}

\begin{ejemplo}[Cartesianas a cilíndricas]
Podemos pasar de coordenadas cartesianas a coordenadas cilíndricas mediante
\begin{equation}
\begin{pmatrix} x^1=x \\ x^2=y \\ x^3=z\end{pmatrix}
\longrightarrow
\begin{pmatrix} x'^1=r \\ x'^2=\phi \\ x'^3=z\end{pmatrix},
\qquad
\begin{aligned}
x&=x(r,\phi)=r\cos\phi,\\
y&=y(r,\phi)=r\sin\phi,\\
z&=z(z)=z .
\end{aligned}
\end{equation}
El vector posición de un punto arbitrario, en ambos sistemas, es
\begin{equation}
\rr(x,y,z) = x\ii + y\jj + z\kk = r\cos\phi\,\ii + r\sin\phi\,\jj + z\kk = \rrp(r,\phi,z).
\end{equation}
En coordenadas cartesianas el vector posición es lineal en $x,y,z$; en coordenadas cilíndricas, en cambio, $\rrp(r,\phi,z)$ es lineal en $r$ y $z$, pero no lo es en $\phi$ (aparece dentro de funciones trigonométricas).
\end{ejemplo}

\begin{ejemplo}[Cartesianas a esféricas]
De manera análoga, podemos pasar de coordenadas cartesianas a esféricas mediante
\begin{equation}
\begin{pmatrix} x^1=x \\ x^2=y \\ x^3=z\end{pmatrix}
\longrightarrow
\begin{pmatrix} x'^1=r \\ x'^2=\phi \\ x'^3=\theta\end{pmatrix},
\qquad
\begin{aligned}
x&=x(r,\phi,\theta)=r\cos\phi\sin\theta,\\
y&=y(r,\phi,\theta)=r\sin\phi\sin\theta,\\
z&=z(r,\theta)=r\cos\theta .
\end{aligned}
\end{equation}
Nuevamente
\begin{equation}
\rr(x,y,z)=x\ii+y\jj+z\kk = r\cos\phi\sin\theta\,\ii + r\sin\phi\sin\theta\,\jj + r\cos\theta\,\kk = \rrp(r,\phi,\theta),
\end{equation}
donde $\rrp(r,\phi,\theta)$ es lineal solo en $r$, y no lo es en $\phi$ y $\theta$.
\end{ejemplo}

\subsection{Coordinatización de \texorpdfstring{$\mathbb{R}^3$}{R3}}

\begin{definicion}
Una \textbf{coordinatización} de $\mathbb{R}^3$ es un mapeo de $M\subseteq\mathbb{R}^3$ en $\mathbb{R}^3$, donde
\begin{equation}
M = I_1\times I_2 \times I_3, \qquad x^1\in I_1,\ x^2\in I_2,\ x^3\in I_3,
\end{equation}
con $I_i\subseteq\mathbb{R}$ ($i=1,2,3$) y $\times$ el producto cartesiano. Así,
\begin{equation}
\rr : M \to \mathbb{R}^3, \qquad x^i \mapsto \rr(x^i).
\end{equation}
\end{definicion}

\begin{ejemplo}
\begin{enumerate}[label=\alph*)]
\item \textbf{Coordenadas cartesianas.} Aquí $x^1,x^2,x^3\in(-\infty,\infty)=\mathbb{R}$, de modo que $M=\mathbb{R}\times\mathbb{R}\times\mathbb{R}=\mathbb{R}^3$ y
\begin{equation}
\rr:\mathbb{R}^3\to\mathbb{R}^3, \qquad (x,y,z)\mapsto \rr(x,y,z).
\end{equation}
\item \textbf{Coordenadas cilíndricas.} Aquí $r\in[0,\infty)$, $\phi\in[0,2\pi)$, $z\in(-\infty,\infty)$, de modo que $M=[0,\infty)\times[0,2\pi)\times\mathbb{R}$ y
\begin{equation}
\rr:M\to\mathbb{R}^3, \qquad (r,\phi,z)\mapsto \rr(r,\phi,z).
\end{equation}
\item \textbf{Coordenadas esféricas.} Aquí $r\in[0,\infty)$, $\phi\in[0,2\pi)$, $\theta\in[0,\pi]$, de modo que $M=[0,\infty)\times[0,2\pi)\times[0,\pi]$ y
\begin{equation}
\rr:M\to\mathbb{R}^3, \qquad (r,\phi,\theta)\mapsto \rr(r,\phi,\theta).
\end{equation}
\end{enumerate}
\end{ejemplo}

\subsection{Invertibilidad: la matriz Jacobiana}

Una transformación de coordenadas es un mapeo de una coordinatización $M$ en otra coordinatización $M'$ de $\mathbb{R}^3$:
\begin{equation}
x'^i : M\to M', \qquad x^j\mapsto x'^i(x^j), \qquad x^i\in M \ \xrightleftharpoons[]{\ \cong\ }\ x'^i\in M'.
\end{equation}
Para que esta transformación sea físicamente útil debe ser \textbf{invertible}, es decir
\begin{equation}
x'^i(x^j) \longleftrightarrow x^i(x'^j).
\end{equation}
Del teorema de la función implícita se sabe que el determinante de la matriz Jacobiana de la transformación,
\begin{equation}
J^i{}_j = \frac{\partial x^i}{\partial x'^j},
\end{equation}
llamado \textbf{Jacobiano} y denotado $J$,
\begin{equation}
J := \det\!\left(J^i{}_j\right) = \det\!\left(\frac{\partial x^i}{\partial x'^j}\right),
\end{equation}
debe ser distinto de cero: $J\neq 0$. La matriz Jacobiana está dada explícitamente por
\begin{equation}
J^i{}_j = \frac{\partial x^i}{\partial x'^j} =
\begin{pmatrix}
\dfrac{\partial x^1}{\partial x'^1} & \dfrac{\partial x^1}{\partial x'^2} & \dfrac{\partial x^1}{\partial x'^3} \\[8pt]
\dfrac{\partial x^2}{\partial x'^1} & \dfrac{\partial x^2}{\partial x'^2} & \dfrac{\partial x^2}{\partial x'^3} \\[8pt]
\dfrac{\partial x^3}{\partial x'^1} & \dfrac{\partial x^3}{\partial x'^2} & \dfrac{\partial x^3}{\partial x'^3}
\end{pmatrix},
\end{equation}
donde el índice $i$ rotula las filas y el índice $j$ rotula las columnas.

\begin{definicion}
Una transformación de coordenadas $x'^i(x^j)$ se dice \textbf{invertible} si el Jacobiano de la transformación es distinto de cero,
\begin{equation}
J := \det\!\left(J^i{}_j\right) = \det\!\left(\frac{\partial x^i}{\partial x'^j}\right) \neq 0 .
\end{equation}
\end{definicion}

\subsection{Ejemplos de cálculo del Jacobiano}

\begin{ejemplo}[Cartesianas a cilíndricas]
Sea la transformación de coordenadas cartesianas a cilíndricas. La matriz Jacobiana es
\begin{equation}
J^i{}_j = \frac{\partial x^i}{\partial x'^j} =
\begin{pmatrix}
\dfrac{\partial x}{\partial r} & \dfrac{\partial x}{\partial \phi} & \dfrac{\partial x}{\partial z} \\[6pt]
\dfrac{\partial y}{\partial r} & \dfrac{\partial y}{\partial \phi} & \dfrac{\partial y}{\partial z} \\[6pt]
\dfrac{\partial z}{\partial r} & \dfrac{\partial z}{\partial \phi} & \dfrac{\partial z}{\partial z}
\end{pmatrix}
=
\begin{pmatrix}
\cos\phi & -r\sin\phi & 0 \\
\sin\phi & r\cos\phi & 0 \\
0 & 0 & 1
\end{pmatrix}.
\end{equation}
El Jacobiano de la transformación es
\begin{equation}
J = \det\!\left(J^i{}_j\right) = r,
\end{equation}
que es distinto de cero si $r\neq 0$. El caso $r=0$ corresponde al eje $z$ ($x=y=0$, $z\in\mathbb{R}$), donde el vector coordenado $\vec e_\phi = \partial \rr/\partial\phi$ se anula. La transformación inversa es
\begin{equation}
r = \sqrt{x^2+y^2}, \qquad \phi = \atan\!\left(\frac{y}{x}\right), \qquad z=z .
\end{equation}
\end{ejemplo}

\begin{ejemplo}[Cartesianas a esféricas]
Sea la transformación de coordenadas cartesianas a esféricas. La matriz Jacobiana es
\begin{equation}
J^i{}_j = \frac{\partial x^i}{\partial x'^j} =
\begin{pmatrix}
\cos\phi\sin\theta & -r\sin\phi\sin\theta & r\cos\phi\cos\theta \\
\sin\phi\sin\theta & r\cos\phi\sin\theta & r\sin\phi\cos\theta \\
\cos\theta & 0 & -r\sin\theta
\end{pmatrix}.
\end{equation}
El Jacobiano de la transformación es
\begin{equation}
J = \det\!\left(J^i{}_j\right) = -r^2\sin\theta,
\end{equation}
que es distinto de cero si $r\neq 0$ y $\theta\neq 0,\pi$. El caso $r=0$ corresponde al origen (un punto) y $\theta=0,\pi$ corresponde al eje $z$ ($x=y=0$, $z\in\mathbb{R}$), donde el vector coordenado $\vec e_\phi = \partial\rr/\partial\phi$ se anula. La transformación inversa es
\begin{equation}
r = \sqrt{x^2+y^2+z^2}, \qquad
\phi = \atan\!\left(\frac{y}{x}\right), \qquad
\theta = \atan\!\left(\frac{\sqrt{x^2+y^2}}{z}\right).
\end{equation}
\end{ejemplo}

\begin{observacion}
La diferencia de signo entre el Jacobiano de la transformación a cilíndricas y el de la transformación a esféricas se debe a que las coordenadas cilíndricas poseen vectores coordenados \textbf{dextrógiros}, mientras que las coordenadas esféricas (en el orden $r,\phi,\theta$ aquí utilizado) poseen vectores coordenados \textbf{levógiros}.
\end{observacion}

\begin{ejemplo}[La Jacobiana inversa en coordenadas cilíndricas]
Calculemos ahora la Jacobiana inversa
\begin{equation}
(J^{-1})^i{}_j = \frac{\partial x'^i}{\partial x^j} =
\begin{pmatrix}
\dfrac{\partial r}{\partial x} & \dfrac{\partial r}{\partial y} & \dfrac{\partial r}{\partial z} \\[6pt]
\dfrac{\partial \phi}{\partial x} & \dfrac{\partial \phi}{\partial y} & \dfrac{\partial \phi}{\partial z} \\[6pt]
\dfrac{\partial z}{\partial x} & \dfrac{\partial z}{\partial y} & \dfrac{\partial z}{\partial z}
\end{pmatrix}
=
\begin{pmatrix}
\dfrac{x}{\sqrt{x^2+y^2}} & \dfrac{y}{\sqrt{x^2+y^2}} & 0 \\[8pt]
-\dfrac{y}{x^2+y^2} & \dfrac{x}{x^2+y^2} & 0 \\[8pt]
0 & 0 & 1
\end{pmatrix}.
\end{equation}
Notemos que $J^i{}_j$ y $(J^{-1})^i{}_j$ están escritas en coordenadas distintas: $J^i{}_j$ es la Jacobiana de la transformación cartesianas $\to$ cilíndricas (las coordenadas cilíndricas son las nuevas), mientras que $(J^{-1})^i{}_j$ es la Jacobiana de la transformación cilíndricas $\to$ cartesianas (las coordenadas cartesianas son las nuevas). Para verificar que efectivamente son matrices inversas, las escribimos en el mismo sistema, por estética en coordenadas cilíndricas:
\begin{equation}
J^i{}_j =
\begin{pmatrix}
\cos\phi & -r\sin\phi & 0 \\
\sin\phi & r\cos\phi & 0 \\
0 & 0 & 1
\end{pmatrix},
\qquad
(J^{-1})^i{}_j =
\begin{pmatrix}
\cos\phi & \sin\phi & 0 \\
-\dfrac{\sin\phi}{r} & \dfrac{\cos\phi}{r} & 0 \\[6pt]
0 & 0 & 1
\end{pmatrix}.
\end{equation}
En efecto,
\begin{equation}
\begin{pmatrix}
\cos\phi & -r\sin\phi & 0 \\
\sin\phi & r\cos\phi & 0 \\
0 & 0 & 1
\end{pmatrix}
\begin{pmatrix}
\cos\phi & \sin\phi & 0 \\
-\dfrac{\sin\phi}{r} & \dfrac{\cos\phi}{r} & 0 \\[6pt]
0 & 0 & 1
\end{pmatrix}
=
\begin{pmatrix}
1 & 0 & 0 \\
0 & 1 & 0 \\
0 & 0 & 1
\end{pmatrix}.
\end{equation}
\end{ejemplo}

% ====================================================================
\section{Vectores base}
% ====================================================================

\subsection{Vectores coordenados}

Cada sistema coordenado $Ox^i$ y $O'x'^i$ posee sus propios \textbf{vectores coordenados},
\begin{equation}
\rr = \rr(x^i), \qquad \ee{i}(x^j) = \frac{\partial \rr(x^j)}{\partial x^i};
\qquad\qquad
\rrp = \rrp(x'^i), \qquad \eep{i}(x'^j) = \frac{\partial \rrp(x'^j)}{\partial x'^i} .
\end{equation}
Un vector arbitrario $\vec A$ se escribe en ambos sistemas como
\begin{equation}
\vec A\Big|_{O} = \sum_{i=1}^{3} A^i(x^j)\,\ee{i}(x^k), \qquad
\vec A\Big|_{O'} = \sum_{i=1}^{3} A'^i(x'^j)\,\eep{i}(x'^k),
\end{equation}
donde $A^i(x^j)$ y $A'^i(x'^j)$ son las componentes de $\vec A$ en los sistemas coordenados $Ox^i$ y $O'x'^j$, respectivamente.

\subsection{Transformación de la base coordenada}

\begin{proposicion}\label{prop:base}
Dados dos conjuntos de vectores coordenados $\{\ee{i}\}_{i=1}^{3}$ y $\{\eep{i}\}_{i=1}^{3}$, asociados a los sistemas coordenados $Ox^i$ y $O'x'^j$, están relacionados bajo la transformación de coordenada $x^i=x^i(x'^j)$ en la forma
\begin{equation}
\eep{i} = \sum_{j=1}^{3} \frac{\partial x^j}{\partial x'^i}\, \ee{j}.
\end{equation}
\end{proposicion}

\begin{proof}
Las coordenadas viejas $Ox^i$ y las nuevas $O'x'^j$ están relacionadas a través del vector posición $\rr\big(x^i(x'^j)\big)=\rrp(x'^j)$. Usando la definición de los vectores base $\{\eep{i}\}_{i=1}^3$ y la regla de la cadena,
\begin{equation}
\eep{i}(x'^j) = \frac{\partial \rrp(x'^j)}{\partial x'^i}
= \frac{\partial \rr\big(x^j(x'^j)\big)}{\partial x'^i}
= \frac{\partial x^1}{\partial x'^i}\frac{\partial \rr}{\partial x^1}
+ \frac{\partial x^2}{\partial x'^i}\frac{\partial \rr}{\partial x^2}
+ \frac{\partial x^3}{\partial x'^i}\frac{\partial \rr}{\partial x^3}
= \sum_{j=1}^{3}\frac{\partial x^j}{\partial x'^i}\frac{\partial \rr}{\partial x^j}
= \sum_{j=1}^{3}\frac{\partial x^j}{\partial x'^i}\, \ee{j}.
\end{equation}
\end{proof}

Para despejar $\ee{i}$ en función de $\eep{i}$ necesitamos la inversa de $\partial x^i/\partial x'^j$. Partiendo de la matriz identidad y usando la regla de la cadena, $x^i=x^i(x'^j)$,
\begin{equation}
\delta^i_j = \frac{\partial x^i}{\partial x^j}
= \sum_{k=1}^{3} \frac{\partial x^i}{\partial x'^k}\frac{\partial x'^k}{\partial x^j},
\end{equation}
y, análogamente, como la transformación inversa $x'^i=x'^i(x^j)$ también debe ser invertible,
\begin{equation}
\delta^i_j = \frac{\partial x'^i}{\partial x'^j} = \sum_{k=1}^{3}\frac{\partial x'^i}{\partial x^k}\frac{\partial x^k}{\partial x'^j}.
\end{equation}
Resumiendo,
\begin{equation}
\sum_{k=1}^{3}\frac{\partial x^i}{\partial x'^k}\frac{\partial x'^k}{\partial x^j}=\delta^i_j,
\qquad
\sum_{k=1}^{3}\frac{\partial x'^i}{\partial x^k}\frac{\partial x^k}{\partial x'^j}=\delta^i_j,
\end{equation}
es decir, la matriz $\partial x'^i/\partial x^j$ es la matriz inversa de $\partial x^i/\partial x'^j$, y viceversa. Si $J^i{}_j = \partial x^i/\partial x'^j$ es la Jacobiana de $x^i=x^i(x'^j)$, entonces su inversa es $(J^{-1})^i{}_j = \partial x'^i/\partial x^j$, de modo que $\sum_{k} J^i{}_k (J^{-1})^k{}_j = \delta^i_j$.

Usando la identidad $\det(AB)=\det(A)\det(B)$ obtenemos, de $\det\!\left(\sum_k J^i{}_k(J^{-1})^k{}_j\right) = \det(\delta^i_j)=1$,
\begin{equation}
\det\!\left(J^i{}_k\right)\det\!\left((J^{-1})^k{}_j\right) = 1
\qquad\Longrightarrow\qquad
\det\!\left((J^{-1})^i{}_j\right) = \frac{1}{J} = J^{-1}.
\end{equation}
Es decir, \textbf{el determinante de la Jacobiana inversa es el inverso del Jacobiano}.

Volviendo a la relación $\eep{i}=\sum_j (\partial x^j/\partial x'^i)\,\ee{j}$ y contrayendo con $\sum_i \partial x'^i/\partial x^k$,
\begin{equation}
\sum_{i=1}^3 \frac{\partial x'^i}{\partial x^k}\,\eep{i}
= \sum_{i=1}^3\sum_{j=1}^3 \frac{\partial x'^i}{\partial x^k}\frac{\partial x^j}{\partial x'^i}\,\ee{j}
= \sum_{j=1}^3\left(\sum_{i=1}^3\frac{\partial x'^i}{\partial x^k}\frac{\partial x^j}{\partial x'^i}\right)\ee{j}
= \sum_{j=1}^3 \delta_{kj}\,\ee{j} = \ee{k},
\end{equation}
de donde se obtiene la ley de transformación inversa de la base coordenada:
\begin{equation}
\ee{k} = \sum_{i=1}^{3} \frac{\partial x'^i}{\partial x^k}\, \eep{i}.
\end{equation}

\begin{observacion}[Resumen]
\begin{itemize}
\item Transformación de coordenadas $x^i\to x'^j$, con $x^i=x^i(x'^j)$: \quad $\displaystyle \eep{i} = \sum_{j=1}^3 \frac{\partial x^j}{\partial x'^i}\,\ee{j}$.
\item Transformación de coordenadas $x'^i\to x^j$, con $x'^i=x'^i(x^j)$: \quad $\displaystyle \ee{i} = \sum_{j=1}^3 \frac{\partial x'^j}{\partial x^i}\,\eep{j}$.
\end{itemize}
\end{observacion}

\subsection{Ejemplo: de cartesianas a esféricas}

Sea la transformación de coordenadas cartesianas a esféricas $x^i=(x,y,z)\leftrightarrow x'^i=(r,\phi,\theta)$. Recordemos que $\ii=(1,0,0)$, $\jj=(0,1,0)$, $\kk=(0,0,1)$. Aplicando $\eep{i}=\sum_j (\partial x^j/\partial x'^i)\,\ee{j}$ para $i=1,2,3$ obtenemos, tras derivar $x=r\cos\phi\sin\theta$, $y=r\sin\phi\sin\theta$, $z=r\cos\theta$ respecto de $r$, $\phi$ y $\theta$:
\begin{align}
\eep{1} &= \vec e_r = \cos\phi\sin\theta\,\ii + \sin\phi\sin\theta\,\jj + \cos\theta\,\kk, \\
\eep{2} &= \vec e_\phi = r\sin\theta\left(-\sin\phi\,\ii + \cos\phi\,\jj\right), \\
\eep{3} &= \vec e_\theta = r\left(\cos\phi\cos\theta\,\ii + \sin\phi\cos\theta\,\jj - \sin\theta\,\kk\right).
\end{align}
Notemos que $\vec e_r$ es un vector unitario, mientras que $\vec e_\phi$ y $\vec e_\theta$ no lo son: $\|\vec e_\phi\|=r\sin\theta$ y $\|\vec e_\theta\|=r$.

Análogamente, efectuando la transformación inversa $x'^i=x'^i(x^j)$ mediante $\ee{i}=\sum_j (\partial x'^j/\partial x^i)\,\eep{j}$ obtenemos
\begin{align}
\ee{1} = \ii &= \frac{x}{\sqrt{x^2+y^2+z^2}}\,\vec e_r - \frac{y}{x^2+y^2}\,\vec e_\phi + \frac{xz}{\sqrt{x^2+y^2}\left(x^2+y^2+z^2\right)}\,\vec e_\theta , \\
\ee{2} = \jj &= \frac{y}{\sqrt{x^2+y^2+z^2}}\,\vec e_r + \frac{x}{x^2+y^2}\,\vec e_\phi + \frac{yz}{\sqrt{x^2+y^2}\left(x^2+y^2+z^2\right)}\,\vec e_\theta , \\
\ee{3} = \kk &= \frac{z}{\sqrt{x^2+y^2+z^2}}\,\vec e_r - \frac{\sqrt{x^2+y^2}}{x^2+y^2+z^2}\,\vec e_\theta .
\end{align}
Por estética y simplicidad, estos mismos vectores se pueden escribir en las propias coordenadas esféricas, usando la base ortonormal $\{\vec e_r,\ \hat e_\phi=\vec e_\phi/(r\sin\theta),\ \hat e_\theta=\vec e_\theta/r\}$:
\begin{align}
\ii &= \cos\phi\sin\theta\,\vec e_r - \frac{\sin\phi}{r\sin\theta}\,\vec e_\phi + \frac{\cos\phi\cos\theta}{r}\,\vec e_\theta, \\
\jj &= \sin\phi\sin\theta\,\vec e_r + \frac{\cos\phi}{r\sin\theta}\,\vec e_\phi + \frac{\sin\phi\cos\theta}{r}\,\vec e_\theta, \\
\kk &= \cos\theta\,\vec e_r - \frac{\sin\theta}{r}\,\vec e_\theta .
\end{align}

% ====================================================================
\section{Transformación de las componentes de un vector}
% ====================================================================

\subsection{Ley de transformación}

\begin{proposicion}\label{prop:comp}
La componente $A^i$ del vector $\vec A$ bajo la transformación de coordenadas $x^i=x^i(x'^j)$ transforma como
\begin{equation}
A'^i = \sum_{j=1}^{3} \frac{\partial x'^i}{\partial x^j}\, A^j .
\end{equation}
\end{proposicion}

\begin{proof}
Como $\vec A$ es un vector, $\vec A\big|_{O'} = \vec A\big|_{O}$, es decir
\begin{equation}
\sum_{i=1}^{3} A'^i \eep{i} = \sum_{j=1}^{3} A^j \ee{j} .
\end{equation}
Usando la ley de transformación de la base, $\ee{j}=\sum_i (\partial x'^i/\partial x^j)\,\eep{i}$,
\begin{equation}
\sum_{i=1}^{3} A'^i \eep{i} = \sum_{j=1}^3\sum_{i=1}^3 A^j \frac{\partial x'^i}{\partial x^j}\eep{i}
= \sum_{i=1}^3\left(\sum_{j=1}^3 \frac{\partial x'^i}{\partial x^j}A^j\right)\eep{i} .
\end{equation}
Luego
\begin{equation}
\sum_{i=1}^{3}\left(A'^i - \sum_{j=1}^3 \frac{\partial x'^i}{\partial x^j}A^j \right)\eep{i} = 0 ,
\end{equation}
y, puesto que la base $\{\eep{i}\}_{i=1}^3$ es linealmente independiente, cada coeficiente debe anularse:
\begin{equation}
A'^i = \sum_{j=1}^{3}\frac{\partial x'^i}{\partial x^j}A^j .
\end{equation}
\end{proof}

Contrayendo esta expresión con $\sum_i \partial x^k/\partial x'^i$ se obtiene la ley de transformación inversa,
\begin{equation}
A^k = \sum_{i=1}^{3}\frac{\partial x^k}{\partial x'^i}\, A'^i .
\end{equation}

\begin{observacion}[Resumen]
\begin{itemize}
\item Transformación $x^i\to x'^j$, con $x^i=x^i(x'^j)$: \quad $\displaystyle A'^i = \sum_{j=1}^3 \frac{\partial x'^i}{\partial x^j}A^j$.
\item Transformación $x'^i\to x^j$, con $x'^i=x'^i(x^j)$: \quad $\displaystyle A^i = \sum_{j=1}^3 \frac{\partial x^i}{\partial x'^j}A'^j$.
\end{itemize}
\end{observacion}

\subsection{Ejemplo: de cartesianas a esféricas}

Un vector $\vec A$ se escribe en coordenadas cartesianas y esféricas como
\begin{equation}
\vec A\Big|_O = \sum_{i=1}^3 A^i\ee{i} = A^x\ii + A^y\jj + A^z\kk, \qquad
\vec A\Big|_{O'} = \sum_{i=1}^3 A'^i\eep{i} = A^r\vec e_r + A^\phi\vec e_\phi + A^\theta\vec e_\theta .
\end{equation}
Aplicando $A'^i=\sum_j(\partial x'^i/\partial x^j)A^j$ con las derivadas $\partial r/\partial x=x/r$, etc., se obtiene, tras pasar a coordenadas esféricas,
\begin{align}
A^r &= \cos\phi\sin\theta\, A^x + \sin\phi\sin\theta\, A^y + \cos\theta\, A^z = \vec e_r\cdot\vec A\Big|_O, \\
A^\phi &= -\frac{\sin\phi}{r\sin\theta}\, A^x + \frac{\cos\phi}{r\sin\theta}\, A^y = \frac{1}{r^2\sin^2\theta}\,\vec e_\phi\cdot\vec A\Big|_O, \\
A^\theta &= \frac{\cos\phi\cos\theta}{r}\, A^x + \frac{\sin\phi\cos\theta}{r}\, A^y - \frac{\sin\theta}{r}\, A^z = \frac{1}{r^2}\,\vec e_\theta\cdot\vec A\Big|_O .
\end{align}
Recíprocamente, aplicando la transformación inversa $A^i=\sum_j (\partial x^i/\partial x'^j)A'^j$ obtenemos
\begin{align}
A^x &= \cos\phi\sin\theta\, A^r - r\sin\phi\sin\theta\, A^\phi + r\cos\phi\cos\theta\, A^\theta
= \ii\cdot \vec A\Big|_{O'}
= \frac{x}{\sqrt{x^2+y^2+z^2}}A^r - y A^\phi + \frac{xz}{\sqrt{x^2+y^2}}A^\theta, \\
A^y &= \sin\phi\sin\theta\, A^r + r\cos\phi\sin\theta\, A^\phi + r\sin\phi\cos\theta\, A^\theta
= \jj\cdot \vec A\Big|_{O'}
= \frac{y}{\sqrt{x^2+y^2+z^2}}A^r + x A^\phi + \frac{yz}{\sqrt{x^2+y^2}}A^\theta, \\
A^z &= \cos\theta\, A^r - r\sin\theta\, A^\phi
= \kk\cdot \vec A\Big|_{O'}
= \frac{z}{\sqrt{x^2+y^2+z^2}}A^r - \sqrt{x^2+y^2}\,A^\theta .
\end{align}

\subsection{Índices arriba y abajo; contracción}

De la discusión anterior se concluye que:
\begin{itemize}
\item las cantidades con el índice \emph{arriba} transforman con la matriz $\partial x'^i/\partial x^j$;
\item las cantidades con el índice \emph{abajo} transforman con la matriz $\partial x^i/\partial x'^j$.
\end{itemize}
En la construcción de un vector cualquiera $\vec A = \sum_i A^i\ee{i}$ se necesitan las cantidades $\{\ee{i}\}_{i=1}^3$ y $\{A^i\}_{i=1}^3$: una tiene el índice abajo y la otra el índice arriba, de modo que se necesita una contracción entre un índice abajo y uno arriba. La razón de esto es que los vectores no dependen del sistema coordenado, como puede verificarse explícitamente: escribiendo $\vec A\big|_{O'} = \sum_i A'^i \eep{i}$ y sustituyendo las leyes de transformación de $A'^i$ y de $\eep{i}$,
\begin{align}
\vec A\Big|_{O'} = \sum_{i=1}^3 A'^i\eep{i}
&= \sum_{i=1}^3\sum_{j=1}^3\frac{\partial x'^i}{\partial x^j}A^j \sum_{k=1}^3\frac{\partial x^k}{\partial x'^i}\ee{k} \notag \\
&= \sum_{j=1}^3\sum_{k=1}^3\left(\sum_{i=1}^3\frac{\partial x'^i}{\partial x^j}\frac{\partial x^k}{\partial x'^i}\right)A^j\ee{k}
= \sum_{j=1}^3\sum_{k=1}^3 \delta_{jk}\,A^j\ee{k} = \sum_{j=1}^3 A^j\ee{j} = \vec A\Big|_{O} ,
\end{align}
como debía ser.

\begin{definicion}
Un índice arriba y un índice abajo se dicen \textbf{contraídos} si están bajo una suma.
\end{definicion}

% ====================================================================
\section{Tensores}
% ====================================================================

\subsection{Tensor contravariante de orden uno}

\begin{definicion}
Una cantidad $T^i$ se dice un \textbf{tensor contravariante de orden uno}, o tensor de orden $(1,0)$, si bajo una transformación de coordenadas $x^i=x^i(x'^j)$ transforma como
\begin{equation}
T'^i = \sum_{j=1}^3 \frac{\partial x'^i}{\partial x^j}\, T^j .
\end{equation}
\end{definicion}
Claramente la componente $A^i$ de un vector arbitrario $\vec A$ es un tensor de orden $(1,0)$.

\subsection{Vectores duales y la base dual}

Los vectores duales también deben satisfacer $f\big|_{O'} = f\big|_{O}$, pues son elementos de un espacio vectorial. De aquí se concluye que las componentes de los vectores duales también deben tener una ley de transformación bajo un cambio de coordenadas. La base dual $\{E^i\}_{i=1}^3$, a diferencia de la base coordenada $\{\ee{i}\}_{i=1}^3$, posee un índice arriba, por lo que las bases duales \emph{no} transforman de la misma manera que las bases coordenadas.

Lo esencial de las bases duales y coordenadas es que la relación que las define,
\begin{equation}
E^i(\ee{j}) = \delta^i_j ,
\end{equation}
debe ser \textbf{invariante}: no depende del sistema coordenado, $E^i(\ee{j})\big|_O = E^i(\ee{j})\big|_{O'} = \delta^i_j$.

\begin{proposicion}\label{prop:dual}
La base dual $\{E^i\}_{i=1}^3$, bajo la transformación de coordenadas $x^i=x^i(x'^j)$, transforma como
\begin{equation}
E'^i = \sum_{j=1}^3 \frac{\partial x'^i}{\partial x^j}\, E^j .
\end{equation}
\end{proposicion}

\begin{proof}
Partimos de $E^i(\ee{j})\big|_{O'} = E'^i(\eep{j})$, y usamos las leyes ya establecidas $\ee{j}=\sum_k (\partial x'^i/\partial x^k)^{-1}\cdots$; explícitamente, escribiendo $\eep{j}$ en función de $\ee{}$ mediante $E^i(\ee{j})$ evaluado con las expansiones apropiadas,
\begin{align}
E^i(\ee{j})\Big|_{O'} = E'^i(\eep{j})
&= \left(\sum_{k=1}^3\frac{\partial x'^i}{\partial x^k}E^k\right)\!\left(\sum_{l=1}^3\frac{\partial x^l}{\partial x'^j}\ee{l}\right) \notag\\
&= \sum_{k=1}^3\sum_{l=1}^3 \frac{\partial x'^i}{\partial x^k}\frac{\partial x^l}{\partial x'^j}\, E^k(\ee{l})
= \sum_{k=1}^3\sum_{l=1}^3 \frac{\partial x'^i}{\partial x^k}\frac{\partial x^l}{\partial x'^j}\,\delta^k_l
= \sum_{k=1}^3 \frac{\partial x'^i}{\partial x^k}\frac{\partial x^k}{\partial x'^j} = \delta^i_j .
\end{align}
Por otro lado esta última cadena de igualdades, leída componente a componente, muestra que
\begin{equation}
E'^i = \sum_{k=1}^3 \frac{\partial x'^i}{\partial x^k}\, E^k .
\end{equation}
\end{proof}

Contrayendo con $\sum_i \partial x^k/\partial x'^i$ se obtiene la ley inversa,
\begin{equation}
E^k = \sum_{i=1}^3 \frac{\partial x^k}{\partial x'^i}\, E'^i .
\end{equation}

\begin{observacion}[Resumen]
\begin{itemize}
\item Transformación $x^i\to x'^j$, con $x^i=x^i(x'^j)$: \quad $\displaystyle E'^i = \sum_{j=1}^3\frac{\partial x'^i}{\partial x^j}E^j$.
\item Transformación $x'^i\to x^j$, con $x'^i=x'^i(x^j)$: \quad $\displaystyle E^i = \sum_{j=1}^3\frac{\partial x^i}{\partial x'^j}E'^j$.
\end{itemize}
\end{observacion}

\subsection{Tensor covariante de orden uno}

\begin{proposicion}
La componente $f_i$ del vector dual $f$ transforma como
\begin{equation}
f'_i = \sum_{j=1}^3 \frac{\partial x^j}{\partial x'^i}\, f_j
\end{equation}
bajo la transformación de coordenadas $x'^i=x'^i(x^j)$.
\end{proposicion}

\begin{proof}
Partimos de $f\big|_O = f\big|_{O'}$, es decir $\sum_j f_j E^j = \sum_i f'_i E'^i$. Usando la ley de transformación de la base dual, $E^j = \sum_i (\partial x^j/\partial x'^i)\,E'^i$,
\begin{equation}
\sum_{j=1}^3 f_j \sum_{i=1}^3 \frac{\partial x^j}{\partial x'^i}E'^i = \sum_{i=1}^3 f'_i E'^i
\quad\Longrightarrow\quad
\sum_{i=1}^3\left(f'_i - \sum_{j=1}^3 \frac{\partial x^j}{\partial x'^i}f_j\right)E'^i = 0 .
\end{equation}
Como $\{E'^i\}_{i=1}^3$ es linealmente independiente,
\begin{equation}
f'_i = \sum_{j=1}^3 \frac{\partial x^j}{\partial x'^i}\, f_j .
\end{equation}
\end{proof}

\begin{definicion}
La cantidad $T_i$ se dice un \textbf{tensor covariante de orden uno}, o tensor de orden $(0,1)$, si bajo una transformación de coordenadas $x^i=x^i(x'^j)$ transforma como
\begin{equation}
T'_i = \sum_{j=1}^3 \frac{\partial x^j}{\partial x'^i}\, T_j .
\end{equation}
\end{definicion}

\begin{observacion}
Es importante aclarar que el término \emph{covarianza} (invariancia de forma) de un vector significa que el vector no cambia de forma bajo una transformación de coordenadas, es decir, no cambia su dirección, sentido ni magnitud. Sin embargo, bajo un cambio de coordenadas cambian las bases coordenadas, por lo que el valor numérico de las componentes del vector sí cambia.
\end{observacion}

% ====================================================================
\section{La métrica}
% ====================================================================

\subsection{Definición y ley de transformación}

La métrica se define mediante el producto punto de los vectores base, $g_{ij}(x^k) = \ee{i}\cdot\ee{j}$.

\begin{proposicion}
La métrica $g_{ij}(x^k)$ bajo una transformación de coordenadas $x^i=x^i(x'^j)$ transforma como
\begin{equation}
g'_{ij} = \sum_{k=1}^3\sum_{l=1}^3 \frac{\partial x^k}{\partial x'^i}\frac{\partial x^l}{\partial x'^j}\, g_{kl} .
\end{equation}
\end{proposicion}

\begin{proof}
Usando la definición de la métrica en coordenadas $O'x'^i$ y la ley de transformación de la base,
\begin{align}
g'_{ij} = \eep{i}\cdot\eep{j}
&= \left(\sum_{k=1}^3\frac{\partial x^k}{\partial x'^i}\ee{k}\right)\cdot\left(\sum_{l=1}^3\frac{\partial x^l}{\partial x'^j}\ee{l}\right) \notag\\
&= \sum_{k=1}^3\sum_{l=1}^3 \frac{\partial x^k}{\partial x'^i}\frac{\partial x^l}{\partial x'^j}\,\ee{k}\cdot\ee{l}
= \sum_{k=1}^3\sum_{l=1}^3 \frac{\partial x^k}{\partial x'^i}\frac{\partial x^l}{\partial x'^j}\, g_{kl} .
\end{align}
\end{proof}

\subsection{Ejemplo: métrica esférica}

Consideremos la transformación de coordenadas cartesianas a esféricas, $(x,y,z)\to(r,\phi,\theta)$. La métrica en coordenadas cartesianas es $g_{ij}=\delta_{ij}$, de modo que
\begin{equation}
g'_{ij} = \sum_{k=1}^3\sum_{l=1}^3 \frac{\partial x^k}{\partial x'^i}\frac{\partial x^l}{\partial x'^j}\,\delta_{kl}
= \sum_{k=1}^3 \frac{\partial x^k}{\partial x'^i}\frac{\partial x^k}{\partial x'^j}
= \frac{\partial x}{\partial x'^i}\frac{\partial x}{\partial x'^j} + \frac{\partial y}{\partial x'^i}\frac{\partial y}{\partial x'^j} + \frac{\partial z}{\partial x'^i}\frac{\partial z}{\partial x'^j} .
\end{equation}
Reemplazando las derivadas de $x=r\cos\phi\sin\theta$, $y=r\sin\phi\sin\theta$, $z=r\cos\theta$ respecto de $(r,\phi,\theta)$ se obtiene
\begin{equation}
g'_{ij} =
\begin{pmatrix}
1 & 0 & 0 \\
0 & r^2\sin^2\theta & 0 \\
0 & 0 & r^2
\end{pmatrix},
\end{equation}
expresión conocida, correspondiente al elemento de línea $\dd s^2 = \dd r^2 + r^2\sin^2\theta\,\dd\phi^2 + r^2\,\dd\theta^2$.

\subsection{Tensor covariante de orden dos}

\begin{definicion}
La cantidad $T_{ij}$ se dice un \textbf{tensor covariante de orden dos}, o tensor de orden $(0,2)$, si transforma como
\begin{equation}
T'_{ij} = \sum_{k=1}^3\sum_{l=1}^3 \frac{\partial x^k}{\partial x'^i}\frac{\partial x^l}{\partial x'^j}\, T_{kl} .
\end{equation}
\end{definicion}
Por lo tanto, la métrica $g_{ij}$ es un tensor de orden $(0,2)$.

% ====================================================================
\section{Invariantes geométricos}
% ====================================================================

\subsection{El elemento de línea}

\begin{proposicion}
El elemento de línea $\dd s$ es invariante bajo transformaciones de coordenadas $x^i=x^i(x'^j)$.
\end{proposicion}

\begin{proof}
Sea el elemento de línea escrito en el sistema $O'x'^i$,
\begin{equation}
\dd s^2\Big|_{O'} = \dd s'^2 = \sum_{i=1}^3\sum_{j=1}^3 g'_{ij}\,\dd x'^i\,\dd x'^j .
\end{equation}
Bajo la transformación $x'^i=x'^i(x^j)$ los diferenciales $\dd x'^i$ transforman como
\begin{equation}
\dd x'^i = \frac{\partial x'^i}{\partial x^1}\dd x^1 + \frac{\partial x'^i}{\partial x^2}\dd x^2 + \frac{\partial x'^i}{\partial x^3}\dd x^3
= \sum_{j=1}^3 \frac{\partial x'^i}{\partial x^j}\dd x^j ,
\end{equation}
de modo que los diferenciales $\dd x'^i$ son tensores de orden $(1,0)$. Sustituyendo esta expresión y la ley de transformación de $g'_{ij}$,
\begin{align}
\dd s^2\Big|_{O'} = \dd s'^2
&= \sum_{i,j,k,l=1}^3 \frac{\partial x^k}{\partial x'^i}\frac{\partial x^l}{\partial x'^j}\, g_{kl}
\left(\sum_{n=1}^3 \frac{\partial x'^i}{\partial x^n}\dd x^n\right)\!\left(\sum_{m=1}^3 \frac{\partial x'^j}{\partial x^m}\dd x^m\right) \notag \\
&= \sum_{k,l,n,m=1}^3\left(\sum_{i=1}^3\frac{\partial x^k}{\partial x'^i}\frac{\partial x'^i}{\partial x^n}\right)\!\left(\sum_{j=1}^3\frac{\partial x^l}{\partial x'^j}\frac{\partial x'^j}{\partial x^m}\right) g_{kl}\,\dd x^n\,\dd x^m \notag\\
&= \sum_{k,l,n,m=1}^3 \delta^k_n\,\delta^l_m\, g_{kl}\,\dd x^n\,\dd x^m
= \sum_{k=1}^3\sum_{l=1}^3 g_{kl}\,\dd x^k\,\dd x^l = \dd s^2\Big|_O .
\end{align}
Por lo tanto $\dd s^2\big|_{O'} = \dd s'^2 = \dd s^2 = \dd s^2\big|_O$, es decir $\dd s = \dd s'$.
\end{proof}

\begin{observacion}
El elemento de línea no es un vector, sino un escalar (un número con dimensión); que sea invariante significa que su valor numérico no cambia. El largo de una trayectoria tendrá el mismo valor en todos los sistemas coordenados.
\end{observacion}

\subsection{El producto escalar}

\begin{proposicion}
El producto escalar entre dos vectores cualesquiera es invariante bajo una transformación de coordenadas.
\end{proposicion}

\begin{proof}
El producto escalar entre dos vectores está dado por $\vec A\cdot\vec B = \sum_{i,j=1}^3 g_{ij}A^i B^j$. Bajo la transformación $x'^i=x'^i(x^j)$,
\begin{equation}
\vec A\cdot\vec B\Big|_{O'} = \sum_{i,j=1}^3 g'_{ij}A'^iB'^j
= \sum_{i,j,k,l=1}^3 \frac{\partial x^k}{\partial x'^i}\frac{\partial x^l}{\partial x'^j}g_{kl}
\left(\sum_{n=1}^3\frac{\partial x'^i}{\partial x^n}A^n\right)\!\left(\sum_{m=1}^3\frac{\partial x'^j}{\partial x^m}B^m\right).
\end{equation}
Reordenando y usando $\sum_i (\partial x^k/\partial x'^i)(\partial x'^i/\partial x^n)=\delta^k_n$ (y análogamente para $l,m$),
\begin{equation}
\vec A\cdot\vec B\Big|_{O'} = \sum_{k,l,n,m=1}^3 \delta^k_n\,\delta^l_m\, g_{kl}A^nB^m
= \sum_{n,m=1}^3 g_{nm}A^nB^m = \vec A\cdot\vec B\Big|_O .
\end{equation}
Por lo tanto $\vec A\cdot\vec B\big|_{O'} = \vec A\cdot\vec B\big|_O$.
\end{proof}

\begin{observacion}
El producto escalar corresponde a un número real, no a un vector. Por lo tanto, su valor numérico no cambia bajo una transformación de coordenadas.
\end{observacion}

\subsection{Magnitud y ángulo entre vectores}

\begin{proposicion}
La magnitud de un vector es invariante bajo transformaciones de coordenadas.
\end{proposicion}

\begin{proof}
Como se demostró, $\vec A\cdot\vec B\big|_{O'} = \vec A\cdot\vec B\big|_O$ para cualesquiera $\vec A,\vec B$. Haciendo $\vec A=\vec B$,
\begin{equation}
\vec A\cdot\vec A\Big|_{O'} = \vec A\cdot\vec A\Big|_O
\qquad\Longrightarrow\qquad
\|\vec A\|^2\Big|_{O'} = \|\vec A\|^2\Big|_O
\qquad\therefore\qquad
\|\vec A\|\Big|_{O'} = \|\vec A\|\Big|_O .
\end{equation}
\end{proof}

Esto significa que el largo de un vector, que es un escalar, no cambia bajo una transformación de coordenadas.

\begin{corolario}
El ángulo entre dos vectores es un invariante bajo una transformación de coordenadas.
\end{corolario}

\begin{proof}
Sabemos que el ángulo $\theta$ entre los vectores $\vec A$ y $\vec B$ está dado por
\begin{equation}
\cos\theta = \frac{\vec A\cdot\vec B}{\|\vec A\|\,\|\vec B\|} .
\end{equation}
Bajo una transformación de coordenadas,
\begin{equation}
\cos\theta\Big|_{O'} = \left.\frac{\vec A\cdot\vec B}{\|\vec A\|\,\|\vec B\|}\right|_{O'}
= \left.\frac{\vec A\cdot\vec B}{\|\vec A\|\,\|\vec B\|}\right|_{O} = \cos\theta\Big|_O ,
\end{equation}
donde se usó la invariancia del producto escalar y de las magnitudes. Por lo tanto $\theta\big|_{O'}=\theta\big|_O$.
\end{proof}

\subsection{Definición formal del producto escalar (dualidad)}

Formalmente, el producto escalar entre un vector dual y un vector está definido como
\begin{equation}
f(v) = \sum_{i=1}^3 f_i\, v^i ,
\end{equation}
donde $v=\sum_i v^i\ee{i}\in V(3,\mathbb{K})$ y $f=\sum_i f_i E^i \in V^*(3,\mathbb{K})$.

\begin{proposicion}
El producto escalar $f(v)=\sum_i f_iv^i$ es invariante bajo la transformación de coordenadas $x'^i=x'^i(x^j)$.
\end{proposicion}

\begin{proof}
Las cantidades $v^i$ y $f_i$ son tensores de orden $(1,0)$ y $(0,1)$, respectivamente, y transforman como
\begin{equation}
v'^i = \sum_{j=1}^3 \frac{\partial x'^i}{\partial x^j}v^j, \qquad
f'_i = \sum_{j=1}^3 \frac{\partial x^j}{\partial x'^i}f_j .
\end{equation}
Entonces
\begin{align}
f'(v') = \sum_{i=1}^3 f'_i v'^i
&= \sum_{i=1}^3\left(\sum_{j=1}^3\frac{\partial x^j}{\partial x'^i}f_j\right)\!\left(\sum_{k=1}^3\frac{\partial x'^i}{\partial x^k}v^k\right) \notag \\
&= \sum_{j,k=1}^3\left(\sum_{i=1}^3\frac{\partial x^j}{\partial x'^i}\frac{\partial x'^i}{\partial x^k}\right)f_jv^k
= \sum_{j,k=1}^3 \delta^j_k\, f_j v^k = \sum_{j=1}^3 f_jv^j = f(v) .
\end{align}
\end{proof}

% ====================================================================
\section{La métrica inversa}
% ====================================================================

\subsection{Relación con la base dual}

\begin{proposicion}
La métrica inversa $g^{ij}$ es dada en términos de la base dual $\{E^i\}_{i=1}^3$ de $\mathbb{R}^3$ en la forma
\begin{equation}
g^{ij} = E^i\cdot E^j .
\end{equation}
\end{proposicion}

\begin{proof}
Sabemos que la base dual está relacionada con la base coordenada a través de la métrica inversa, $E^i = \sum_{j=1}^3 g^{ij}\ee{j}$, de modo que
\begin{equation}
E^i(\ee{j}) = E^i\cdot\ee{j} = \sum_{k=1}^3 g^{ik}\,\ee{k}\cdot\ee{j} = \sum_{k=1}^3 g^{ik}g_{kj} = \delta^i_j .
\end{equation}
Luego,
\begin{align}
E^i\cdot E^j = \left(\sum_{k=1}^3 g^{ik}\ee{k}\right)\cdot\left(\sum_{l=1}^3 g^{jl}\ee{l}\right)
&= \sum_{k=1}^3\sum_{l=1}^3 g^{ik}g^{jl}\,\ee{k}\cdot\ee{l} = \sum_{k=1}^3\sum_{l=1}^3 g^{ik}g^{jl}g_{kl} \notag \\
&= \sum_{k=1}^3 g^{ik}\left(\sum_{l=1}^3 g^{jl}g_{kl}\right) = \sum_{k=1}^3 g^{ik}\delta^j_k = g^{ij} .
\end{align}
\end{proof}

\subsection{Ejemplo: métrica inversa esférica}

Calculemos la métrica inversa asociada a las coordenadas esféricas usando $g^{ij}=E^i\cdot E^j$, donde $E^1=E^r$, $E^2=E^\phi$, $E^3=E^\theta$ son la base dual asociada a $\vec e_r$, $\vec e_\phi=r\sin\theta\,\hat e_\phi$, $\vec e_\theta = r\,\hat e_\theta$. Puesto que la base coordenada esférica es ortogonal, la base dual coincide con $E^1=\vec e_r$, $E^2=\vec e_\phi/(r^2\sin^2\theta)$, $E^3=\vec e_\theta/r^2$, de donde
\begin{align}
g^{11} &= E^1\cdot E^1 = \vec e_r\cdot\vec e_r = 1, \\
g^{12}=g^{21} &= E^1\cdot E^2 = \vec e_r\cdot \frac{\vec e_\phi}{r^2\sin^2\theta} = 0, \\
g^{13}=g^{31} &= E^1\cdot E^3 = \vec e_r\cdot\frac{\vec e_\theta}{r^2} = 0, \\
g^{22} &= E^2\cdot E^2 = \frac{\vec e_\phi}{r^2\sin^2\theta}\cdot\frac{\vec e_\phi}{r^2\sin^2\theta}
= \frac{r^2\sin^2\theta}{r^4\sin^4\theta} = \frac{1}{r^2\sin^2\theta}, \\
g^{23}=g^{32} &= E^2\cdot E^3 = \frac{\vec e_\phi}{r^2\sin^2\theta}\cdot\frac{\vec e_\theta}{r^2} = 0, \\
g^{33} &= E^3\cdot E^3 = \frac{\vec e_\theta}{r^2}\cdot\frac{\vec e_\theta}{r^2} = \frac{r^2}{r^4} = \frac{1}{r^2} .
\end{align}
Resumiendo,
\begin{equation}
g^{ij} =
\begin{pmatrix}
1 & 0 & 0 \\
0 & \dfrac{1}{r^2\sin^2\theta} & 0 \\[6pt]
0 & 0 & \dfrac{1}{r^2}
\end{pmatrix},
\end{equation}
resultado conocido.

\subsection{Ley de transformación y tensor contravariante de orden dos}

\begin{proposicion}
La métrica inversa $g^{ij}$ bajo una transformación de coordenadas $x'^i=x'^i(x^j)$ transforma como
\begin{equation}
g'^{\,ij} = \sum_{k=1}^3\sum_{l=1}^3 \frac{\partial x'^i}{\partial x^k}\frac{\partial x'^j}{\partial x^l}\, g^{kl} .
\end{equation}
\end{proposicion}

\begin{proof}
Escribiendo la métrica inversa en coordenadas $O'x'^i$, $g'^{\,ij} = E'^i\cdot E'^j$, y usando la ley de transformación de la base dual, $E'^i = \sum_k(\partial x'^i/\partial x^k)E^k$,
\begin{align}
g'^{\,ij} = E'^i\cdot E'^j
&= \left(\sum_{k=1}^3\frac{\partial x'^i}{\partial x^k}E^k\right)\cdot\left(\sum_{l=1}^3\frac{\partial x'^j}{\partial x^l}E^l\right) \notag \\
&= \sum_{k=1}^3\sum_{l=1}^3 \frac{\partial x'^i}{\partial x^k}\frac{\partial x'^j}{\partial x^l}\, E^k\cdot E^l
= \sum_{k=1}^3\sum_{l=1}^3 \frac{\partial x'^i}{\partial x^k}\frac{\partial x'^j}{\partial x^l}\, g^{kl} .
\end{align}
\end{proof}

\begin{definicion}
La cantidad $T^{ij}$ se dice un \textbf{tensor contravariante de orden dos}, o tensor de orden $(2,0)$, si bajo una transformación de coordenadas $x'^i=x'^i(x^j)$ transforma como
\begin{equation}
T'^{ij} = \sum_{k=1}^3\sum_{l=1}^3 \frac{\partial x'^i}{\partial x^k}\frac{\partial x'^j}{\partial x^l}\, T^{kl} .
\end{equation}
\end{definicion}
Por lo tanto, la métrica inversa $g^{ij}$ es un tensor de orden $(2,0)$.

\subsection{Ejemplo: verificación directa}

Calculemos la métrica inversa asociada a las coordenadas esféricas directamente a partir de la métrica inversa asociada a las coordenadas cartesianas, $g^{kl}=\delta^{kl}$:
\begin{equation}
g'^{\,ij} = \sum_{k=1}^3\sum_{l=1}^3\frac{\partial x'^i}{\partial x^k}\frac{\partial x'^j}{\partial x^l}\,\delta^{kl}
= \frac{\partial x'^i}{\partial x}\frac{\partial x'^j}{\partial x} + \frac{\partial x'^i}{\partial y}\frac{\partial x'^j}{\partial y} + \frac{\partial x'^i}{\partial z}\frac{\partial x'^j}{\partial z} .
\end{equation}
Realizaremos el cálculo para $i=j=1$ (es decir, $g'^{11}$); los demás se obtienen de forma análoga:
\begin{align}
g'^{11} &= \left(\frac{\partial r}{\partial x}\right)^{\!2} + \left(\frac{\partial r}{\partial y}\right)^{\!2} + \left(\frac{\partial r}{\partial z}\right)^{\!2} \notag \\
&= \left(\frac{\partial \sqrt{x^2+y^2+z^2}}{\partial x}\right)^{\!2} + \left(\frac{\partial \sqrt{x^2+y^2+z^2}}{\partial y}\right)^{\!2} + \left(\frac{\partial \sqrt{x^2+y^2+z^2}}{\partial z}\right)^{\!2} \notag \\
&= \frac{x^2}{r^2} + \frac{y^2}{r^2} + \frac{z^2}{r^2} = 1 .
\end{align}
Efectuando los cálculos restantes se obtiene, nuevamente,
\begin{equation}
g^{ij} =
\begin{pmatrix}
1 & 0 & 0 \\
0 & \dfrac{1}{r^2\sin^2\theta} & 0 \\[6pt]
0 & 0 & \dfrac{1}{r^2}
\end{pmatrix},
\end{equation}
resultado conocido, consistente con el obtenido en la sección anterior.

% ====================================================================
\section{Las deltas de Kronecker}
% ====================================================================

La métrica inversa $g^{ij}$ es un tensor de orden $(2,0)$. Recordemos que existen tres cantidades que aparentemente representan la matriz identidad: $\delta_{ij}$, $\delta^{ij}$ y $\delta^i_j$. Sin embargo, solo una de ellas representa realmente la identidad: $\delta^i_j$. En efecto, esta cantidad es invariante bajo transformaciones de coordenadas:
\begin{equation}
\delta'^{\,i}_j = \sum_{k=1}^3\sum_{l=1}^3 \frac{\partial x'^i}{\partial x^k}\frac{\partial x^l}{\partial x'^j}\,\delta^k_l
= \sum_{k=1}^3 \frac{\partial x'^i}{\partial x^k}\frac{\partial x^k}{\partial x'^j}
= \frac{\partial x'^i}{\partial x'^j} = \delta^i_j .
\end{equation}

En cambio, las cantidades $\delta_{ij}$ y $\delta^{ij}$ transforman como tensores de orden $(0,2)$ y $(2,0)$ respectivamente, y \emph{no} son invariantes, pues cambian su forma funcional; por ejemplo, al pasar a coordenadas esféricas dejan de ser la matriz identidad (como vimos, se convierten en $\mathrm{diag}(1,r^2\sin^2\theta,r^2)$ y $\mathrm{diag}(1,1/(r^2\sin^2\theta),1/r^2)$, respectivamente).

En conclusión:
\begin{itemize}
\item La cantidad $\delta^i_j$ es \textbf{la matriz identidad} (invariante en todo sistema coordenado).
\item La cantidad $\delta_{ij}$ es la métrica asociada a las coordenadas cartesianas, llamada \textbf{métrica de Euclides}.
\item La cantidad $\delta^{ij}$ es la métrica inversa de la métrica de Euclides.
\end{itemize}
